\section{Where the time goes}
\label{sec:profile}

I profiled with \texttt{sample}, the macOS profiler, which records the running
instruction once per millisecond; a run of \inst{php11\_10} gives about
3{,}400 samples, so a share is accurate to about $\pm 1$ percentage point. On
that instance \texttt{propagate()} takes \numPropagateShare\% of the run and
\texttt{analyze()} \numAnalyzeShare\%; everything else is below 1\%. Cheap is
not the same as unimportant, though: clause deletion takes almost no time
itself, yet it made \inst{php11\_10} 4.4 times faster by keeping the clause
database small.

Inside \texttt{propagate()}, about half of the whole run goes into searching a
clause for a new literal to watch after one of its watched literals has become
false (Figure~\ref{fig:regions}), and one instruction alone receives 37\% of
all samples. A sample lands on the instruction the processor is \emph{waiting
at}, not on the one that caused the wait, an effect called skid. The hot
instruction is a cheap load from L1, reachable only by a jump from the check
that ends the search loop. That check depends on the current assignment and is
hard to predict, and after a wrong guess the processor restarts exactly there.
This points to branch mispredictions, a hypothesis I cannot confirm without
hardware counters.

The random instance supports it. Its learnt clauses are shorter, so a search
checks about 2 literals instead of 10: the loop body shrinks from 8\% to
1.5\% of the run, but the loop exit stays expensive, so the cost is paid per
search, not per literal. Moving a watch, a store to an essentially random
address, costs only 2.4\%, because the processor waits for loads but not for
stores.

\begin{figure}[ht]
\centering
\plotfigure{figures/fig-regions}
\caption{Share of all profile samples in each part of \texttt{propagate()}.
Long learnt clauses (pigeonhole) make the search for a new watch dominant;
short ones (random 3-SAT) shift weight to the fixed cost of each visit.}
\label{fig:regions}
\end{figure}
